\section{Concise proofs for the anchored connection}
\label{sec:lean-anchor-proofs}

This section proves the results of
\cref{sec:anchored-physical-connection}.  The only outer inputs are the
critical splitting, exact quotient, coherent reference, and generated-core
estimates in
\cref{prop:anchor-toolkit-critical-splitting,%
prop:anchor-toolkit-resource-phase-quotient,%
prop:anchor-toolkit-reference-hits,prop:anchor-toolkit-generated-core}.
In particular, we do not repeat
the original-endpoint calculation.  Instead, both physical histories are
compared, in one common complete-history chart, with the two histories of
the finite selected problem.  The finite-gap estimate of
\cref{prop:finite-gap-bridge} can then be transferred directly to the
intrinsic physical gap.

Throughout this section, \(M_{\rm loss},R_b,Q_*,p_o\), and
\(\mathbf p_{\rm sel}\) are chosen once, in that order, as in
\cref{prop:anchor-central-finite-shooting}; \(p_c\) is the central Gaussian
exponent of this one preparation.  Every selected object below---the attracting trace, the repelling
trace \(C_{b,p}\), the finite graph, and
\(\mathscr D_N^{\rm sel}\)---comes from this same datum.  The auxiliary
outer exponent \(p_o=(\chi_b/c_b)^2p_c\) fixes only the independently
generated physical reference; it is not a second preparation.  We also use
\[
 \mathscr Y_N=C([-\theta_m,0];\mathbb R^N)\times\mathbb R
\]
for a fold-time voltage history and its current resource value.  The exact
anchored fold-time field used below is
\[
 x_\phi:=\phi(0)-\mathbf1.
\]
\begin{align*}
 \mathscr V^{\rm f}_{N,\delta,p}(\phi,w)
 &=
 \delta^{-1}\!\left[
  \frac23\mathbf1-w\mathbf1
  -c_N\circ x_\phi^{\circ2}
  -\frac\beta3x_\phi^{\circ3}
  +D(P_N-\Id)\phi(0)\right]\\
 &\quad
  +\delta K\sum_{k=0}^mL_{k,N}(\eta)
       \{\phi(0)-\phi(-\theta_k)\},\\
 \mathscr F^{\rm f}_{N,\delta,p}(\phi,w)
 &=\left(\mathscr V^{\rm f}_{N,\delta,p}(\phi,w),
  \delta\,\sigma_{\rm anc}
   \bigl(\pi_N^Tx_\phi;\delta,\nu\bigr)\right).
\end{align*}
The symbol \(\mathcal I_p\) denotes the exact invariant-history lift from
\cref{lem:shared-graph-trace}, followed by reconstruction in this physical
fold-time history chart.
Every \(D_\eta\) and \(J_p^1\) below is a Fr\'echet derivative on
\(\mathfrak T_N^{\rm red}\), measured in the full operator norm; the
selected solution jets are obtained first on affine parameter lines and
then polarized, as in the structural-ball finite-gap theorem.

\begin{proof}[Proof of
  \cref{prop:anchor-central-identity-equilibria}]
On the depth-two retained tube one has \(g_{\rm anc}=1\), so the anchored
and original vector fields agree there as maps of complete histories.  At a
constant history all current-minus-delay terms vanish, while
\cref{eq:anchor-toolkit-critical-transverse,%
eq:anchor-toolkit-critical-curve} annihilate the
remaining voltage field.  Since \(g_{\rm anc}(\pm r_A)=0\), the histories
in \cref{eq:anchor-equilibria} are equilibria.  Finally,
\[
 \partial_r\sigma_{\rm anc}(\pm r_A;\delta,\nu)
   =(\pm r_A-\delta^2\nu)g_{\rm anc}'(\pm r_A)<0,
\]
which proves \cref{eq:anchor-source-derivative-sign}.
\end{proof}

\subsection{The two anchored history objects}

We first extract from the outer toolkit the invariant objects used below.
The short characteristic calculation is included to make the two Morse
indices transparent.

\begin{lemma}[Anchored branches and the intrinsic terminal sheet]
\label[lemma]{lem:lean-anchor-history-objects}
Uniformly over the bounded anchor and structural classes, the following
objects exist after the ordered decrease in
\cref{prop:anchor-toolkit-critical-splitting,%
prop:anchor-toolkit-resource-phase-quotient,%
prop:anchor-toolkit-reference-hits,prop:anchor-toolkit-generated-core}.

\begin{enumerate}[label=\textup{(\roman*)}]
\item Each equilibrium \(E_N^\pm\) has exactly one
open-right-half-plane characteristic root.  At \(E_N^+\) it is the unique
small slow root in \cref{eq:anchor-slow-roots}; at \(E_N^-\) that slow root
is stable and the unique unstable root is the simple strong voltage root.  Thus
\(\dim W^u(E_N^+)=1\) and \(\operatorname{codim}W^s(E_N^-)=1\).
The local invariant manifolds are \(C^2\) on
\(\rho_{A,\delta}=c_A\delta^4\), with
\(d_{A,\delta}=\rho_{A,\delta}/2\); in the scaled charts their state norms
and selected first parameter jets are at most \(C\delta^{-M_A}\).

\item The inward component of \(W^u(E_N^+)\) has a unique past-complete
orbit.  Evolving it only forward gives a unique hit \(A^0_p\) of
\(\Sigma_0\), with a selected first \(p=(\nu,\eta)\)-jet.  On every fixed
positive outer subannulus it is attracted in the event-normal fading norm.

\item There is an exact orbit \(\widehat Z_-(s;p)\), starting on a fixed
regular seed section and converging to \(E_N^-\).  After more than three
delay buffers, recut it at the strong entry section and call the resulting
generated orbit \(\bar Z_-\).  On a fixed regular negative section \(L\), the
histories converging to \(E_N^-\) form a codimension-one \(C^2\) graph with
selected first \(p\)-jet on a ball of radius
\(R_{{\rm hl},\delta}=c_{\rm hl}\delta^6\).  Its full compatible-history membership function
\(\mathfrak m_{L,p}\), at the graph center \(\bar Z_{-,L}\), has a
normalized differential
\[
 \Lambda_{L,p}^{\infty}
 :=D\mathfrak m_{L,p}(\bar Z_{-,L}),
 \qquad
 \Lambda_{L,p}^{\infty}\jmath_{L,p}^{u,\infty}=1.
\]
Only the transported one-dimensional unstable line is inverted in the
half-line construction.  If the common finite core is cut after a buffer
\(L_\delta\), its capped stable/unstable action is at least
\(c\log(1/\delta)/\delta\).

\item If an exact outer reference is recut at \(L\), its event-normal
unstable column \(h_{L,p}^{u,o}\) lies in the same full tangent space and
\begin{align}
 \|J_p^1(h_{L,p}^{u,o}-\jmath_{L,p}^{u,\infty})\|_{\mathscr H_L^2}
 &\le C\delta^{-M}e^{-a/\delta},\notag\\
 \|J_p^1\{D\mathfrak m_{L,p}(O_{L,p})
              -\Lambda_{L,p}^{\infty}\}\|_{(\mathscr H_L^2)^*}
 &\le C\delta^{-M}e^{-a/\delta}.
 \label{eq:lean-anchor-terminal-history-comparison}
\end{align}
Consequently
\begin{equation}
 D\mathfrak m_{L,p}(O_{L,p})h_{L,p}^{u,o}\ge\frac34 .
 \label{eq:lean-anchor-terminal-crossing}
\end{equation}
The membership value at \(O_{L,p}\), with its selected first parameter
jet, has the same bound.

\item Recut the attracting trace of \(\mathbf p_{\rm sel}\), after
literal agreement has begun, on the same positive regular section as the
physical incoming branch.  Write its hit of \(\Sigma_0\) as
\(\mathfrak C_{0,p}(\bar\psi_p,u_p^a)\), and write the physical hit as
\(\mathfrak C_{0,p}(\psi_A,u_A)\).  After the common attracting buffer,
\begin{equation}
 \bigl\|J_p^1\{(\psi_A,u_A)
       -(\bar\psi_p,u_p^a)\}\bigr\|
 \le C\delta^{Q_*+M_{\rm loss}+4},
 \qquad J_p^1\in\{1,\partial_\nu,D_\eta\}.
 \label{eq:lean-anchor-positive-selected-flatness}
\end{equation}
The exponent is fixed before \(\delta_0\) is decreased.
\end{enumerate}
\end{lemma}

\begin{proof}
Put
\[
 J_\pm=\mathcal J_N^{\rm fv}(\pm r_A),\qquad
 A_\pm(\lambda)=J_\pm+\delta^2K\sum_{k=0}^mL_{k,N}(\eta)
       \{1-e^{-\lambda\theta_k/\delta}\}.
\]
At either anchor the modal equations are
\begin{align*}
 (\lambda\Id-A_\pm(\lambda))u&=-\omega\mathbf1,\\
 \lambda\omega&=\delta^2\sigma_\pm'\pi_N^Tu.
\end{align*}
Eliminating \(u\) near the origin gives
\[
 \Delta_\pm(\lambda)
 =\lambda+\delta^2\sigma_\pm'\pi_N^T
      (\lambda\Id-A_\pm(\lambda))^{-1}\mathbf1.
\]
The differentiated critical-curve identity implies
\(\pi_N^T(-J_\pm)^{-1}\mathbf1=-1/w'_{0,N}(\pm r_A)\).
The uniform resolvent estimate in
\cref{prop:anchor-toolkit-critical-splitting} therefore yields
\[
 \Delta_\pm(\lambda)
 =\lambda-\delta^2\frac{\sigma_\pm'}{w'_{0,N}(\pm r_A)}
       +O(\delta^4),\qquad \Delta_\pm'=1+O(\delta^2).
\]
Rouch\'{e}'s theorem gives the small simple root.  The bordered homotopy
which switches the second row from zero to
\(\delta^2\sigma_\pm'\pi_N^T u\) is made explicit by
\[
 \Delta_{\pm,\varepsilon}(\lambda)
 =\lambda+\varepsilon\sigma_\pm'\pi_N^T
  (\lambda\Id-A_\pm(\lambda))^{-1}\mathbf1,
 \qquad 0\le\varepsilon\le\delta^2.
\]
The frozen current dichotomy and the \(O(\delta^2)\) delay pencil give
\[
 \sup_{\xi\in\mathbb R}
 \|(i\xi\Id-A_\pm(i\xi))^{-1}\|\le C.
\]
On a fixed pole-free imaginary segment,
\[
 \Delta_{\pm,\varepsilon}(i\xi)
 =i\xi-\varepsilon
   \frac{\sigma_\pm'}{w'_{0,N}(\pm r_A)}
   +\varepsilon R_\pm(i\xi),\qquad
 |R_\pm(i\xi)|\le C|\xi|.
\]
The real term excludes zeros for
\(|\xi|\le C_0\varepsilon\); after \(C_0\) is chosen large, \(i\xi\)
dominates for \(C_0\varepsilon\le|\xi|\le r_*\); and the uniform
imaginary-axis resolvent excludes them for \(|\xi|\ge r_*\).
Moreover,
\[
 \det(\lambda\Id-A_\pm(\lambda))
       \Delta_{\pm,\varepsilon}(\lambda)
\]
is the entire bordered characteristic determinant, and the same resolvent
estimate on a common right half disk prevents roots from escaping through
infinity.  Removing the small half disk containing the moving root and
applying the argument principle therefore preserves every other
right-half-plane root count along the homotopy.  Hence all other roots have
the frozen voltage count.  This proves
the first sentence of (i); the signs follow from
\cref{eq:anchor-toolkit-frozen-identities,%
eq:anchor-source-derivative-sign}.

The frozen history splitting in
\cref{prop:anchor-toolkit-critical-splitting}, the bordered characteristic
count above, and exponential-trichotomy roughness give the local trichotomy
at each hyperbolic anchor.  The resource shear in
\cref{prop:anchor-toolkit-resource-phase-quotient} is used only on fixed
regular subannuli separated from the zeros of \(q\); it is not used at an
anchor, where \(q^{-1}\) is unbounded.  The scaled local
Lyapunov--Perron contraction is uniform after
the slow coordinate is divided by \(\rho_{A,\delta}\), proving the rest of
(i).  Let \(\rho_+\) be the signed slow coordinate at \(E_N^+\), \(Q_+\) the
fast complementary projection, and \(\Gamma_+\) its finite slow expansion.
On the local collar,
\[
 \rho_+'=\delta\rho_+\{\chi_++O(\rho_+)+O(\delta)\},\quad \chi_+\ge c>0,
 \qquad
 \|Q_+(Z_+-Z_{+,{\rm crit}})\|_\sharp\le C\delta^2|\rho_+|.
\]
Continue the exact inward local-manifold hit to a fixed regular section using
\(\Gamma_+\) and the causal
event-normal Green operator.  The local-to-regular continuation and its
selected first jet have only the fixed loss \(\delta^{-2}\), and the
component ledger is
\begin{equation}
 \|R_{\Gamma_+}\|+\|D_pR_{\Gamma_+}\|\le C\delta^5,\qquad
 \|\mathbb G_+\|\le C\delta^{-2},\qquad
 C\delta^{-2}\|R_{\Gamma_+}\|\le\tfrac14c_+\delta^2,\qquad
 \operatorname{Lip}\{\mathbb G_+\mathcal N_{\Gamma_+}\}\le\tfrac13.
 \label{eq:lean-anchor-positive-ledger}
\end{equation}
The contraction is taken on the radius-\(c_+\delta^2\) component ball.
On the first fixed subannulus the fast displacement is \(O(\delta^2)\), the
delay drift is \(O(\delta)\), and the reference event speed is at least
\(c\delta\), with an \(O(\delta^2)\) correction.  Thus first exit through a
side boundary is impossible, and the continuation overlaps the regular
resource quotient with one transverse hit.
After three generated delay buffers, forward uniqueness and the event
implicit-function theorem give (ii), without a backward semiflow.

For the negative side, choose a regular seed section between the
fold and \(E_N^-\).  Successive solution of the frozen voltage-resource
border and the scalar recovery row gives a finite slow expansion
\(\Gamma_-\) whose literal-history residual, including the old-history
interval, is \(O(\delta^9)\) with one selected parameter derivative.  More
explicitly, on the fixed seed annulus the border
\[
 (Z,Q)\longmapsto
 \mathcal J_N^{\rm fv}(r)Z-Q\mathbf1,\qquad
 (Z,Q)\in E_N\times\mathbb R,
\]
has a uniform inverse.  Starting with
\(q_0(r)=r g_{\rm anc}(r)/w'_{0,N}(r)\), define the delayed feet by the
flow of the already known truncation of \(q\).  At order \(j\), coefficient
extraction in the voltage equation determines \((Z_j,Q_j)\) through this
border, and the resource equation then determines \(q_j\) through the
nonzero factor \(w'_{0,N}(r)\).  Ten steps leave the displayed residual;
Taylor's formula applies equally on the prescribed old-history interval
because every delayed foot stays in the fixed annulus.  On the part ending at
the regular seed section, the generalized resource-phase quotient retains
the factor \(\sigma_{\rm anc}'\).  Concatenating its dichotomy with the local
anchor trichotomy from (i), whose bundles agree there by forward uniqueness,
across their common regular section gives a
half-line mixed Green operator.  In the fading solution and
\(\delta^{-1}\)-scaled range norms, put
\(R_{{\rm hl},\delta}=c_{\rm hl}\delta^6\).  The quantitative ledger is
\begin{align}
 \|\mathbb G_-\|&\le C_G\delta^{-2},&
 \|R_{\Gamma_-}\|+\|D_pR_{\Gamma_-}\|&\le C_R\delta^9,\notag\\
 C_G\delta^{-2}C_R\delta^9
 &\le\tfrac14R_{{\rm hl},\delta},&
 \sup_{\|z\|,\|y\|\le R_{{\rm hl},\delta}}
 \frac{\|\mathbb G_-\{\mathcal N_{\Gamma_-}(z)
             -\mathcal N_{\Gamma_-}(y)\}\|}
      {\|z-y\|}
 &\le Cc_{\rm hl}\delta^{6-2-2}\le\tfrac13 .
 \label{eq:lean-anchor-negative-ledger}
\end{align}
It propagates the stable complement forward
and inverts only the one-dimensional unstable coefficient.  Therefore
\[
 z=\mathbb G_-\bigl\{-R_{\Gamma_-}+\mathcal N_{\Gamma_-}(z)\bigr\}
\]
is a contraction on a radius-\(c\delta^6\) ball.  Its zero gives the exact
seed orbit \(\widehat Z_-\); recutting it after the generated buffers gives
\(\bar Z_-\).  Allowing the stable entry coordinate to vary in the
same fixed-point equation gives the intrinsic terminal graph in (iii).
Differentiating this selected fixed-point equation proves its stated jets.

For (iv), put the exact outer reference and \(\widehat Z_-\) in a common
regular \(r\)-clock between two fixed negative sections.  Thus the time
tangent is quotiented before subtraction and no separate moving-event term
occurs.  Their difference
satisfies the two-sided dichotomy identity from the toolkit.  The stable
left and unstable right actions are at least \(a/\delta\), so absorption
gives \(C\delta^{-M}e^{-a/\delta}\), also after one selected parameter
derivative.  Apply the same estimate to the event-normal columns and use
the \(C^2\) membership row in the common full-history chart.  Taylor's
formula proves \cref{eq:lean-anchor-terminal-history-comparison}; use the
displayed normalization at the terminal sheet.  The resulting scalar
is \(1+o(1)\), which proves
\cref{eq:lean-anchor-terminal-crossing}.

For (v), compare the two exact positive histories only after both have
entered the same anchored equation and freeze their common event clock.  In
the event-normal quotient their difference satisfies the causal
variation-of-constants identity.  The initial discrepancy and its selected
first jet are polynomially bounded, while the fixed attracting buffer has
capped action at least \(c\log(1/\delta)/\delta\).  The nonlinear Volterra
term is absorbed by the toolkit's common contraction margin.  Continuing
through the exact central graph and using its selected trace estimate gives
\(C\delta^{-M}e^{-c\log(1/\delta)/\delta}\), plus the Gaussian collar tail
\(C\delta^{-M}e^{-S_c^2/2+cS_c}\).  The ordered inequalities for \(p_c\)
make both terms at most \(C\delta^{Q_*+M_{\rm loss}+4}\), proving
\cref{eq:lean-anchor-positive-selected-flatness}.
\end{proof}

\begin{proof}[Proof of \cref{thm:anchor-indices-manifolds}]
The characteristic count and the scaled local charts are
\cref{lem:lean-anchor-history-objects}(i).  The local unstable-manifold
chart gives the unique past-complete extension of its inward branch.  The
stable manifold is defined by forward convergence, so it requires no
ambient backward RFDE evolution.  Event transversality and the selected
jet bounds are part of the same local toolkit.  This proves every assertion
of the theorem, including the stated distinction between the later
half-line chart center and an invariant one-dimensional stable family.
\end{proof}

The central passage needs one ambient statement in addition to the
graph-tangent trace inverse.  We record it here because it prevents a
graph inverse from being used as an inverse for arbitrary RFDE histories.

\begin{lemma}[Causal normal inverse in the central chart]
\label[lemma]{lem:lean-anchor-central-normal-inverse}
Let \(Z^c_{p,s}=\mathcal I_p(u^c_p(s))\) be a retained exact graph member
on a fixed-clock central segment whose complete delay enlargement stays in
the exact agreement tube.  With \(e_{\rm pl}\) denoting present planar
evaluation, the shear
\begin{equation}
 \phi\longmapsto
 \bigl(e_{\rm pl}\phi,
       \phi-\mathcal I_p(e_{\rm pl}\phi)\bigr)=:(u,\omega)
 \label{eq:lean-anchor-central-shear}
\end{equation}
is a local \(C^2\)-in-state, \(C^1\)-in-\(p\) diffeomorphism.  On the
normal fiber \(\ker e_{\rm pl}\), the linear path residual together with
the initial compatible history is a causal isomorphism.  Its inverse and
the solutions obtained after one selected \(p\)-derivative have norm at
most \(C\delta^{-M_\perp}\), uniformly in \(N\).
\end{lemma}

\begin{proof}
The exact lift satisfies \(e_{\rm pl}\mathcal I_p=\Id\), so the inverse of
\cref{eq:lean-anchor-central-shear} is
\((u,\omega)\mapsto\mathcal I_p(u)+\omega\).  Subtracting the derivative of
the graph invariance equation from the full variational equation makes the
normal block triangular.  Its planar-memory component is a left shift plus
an integral over at most one maximal delay, and its transverse current
component is
\[
 \rho'(s)=\delta^{-1}A_N\rho(s)
          +\mathscr C_p(s)[\kappa_s,\rho_s]+f(s).
\]
The Poisson estimate gives
\[
 \left\|\int_0^s e^{A_N(s-t)/\delta}f(t)\,\dd t\right\|_{C(E_N)}
 \le \frac{\delta}{D\gamma}\|f\|_{C(E_N)}.
\]
The feedback from the normal history to the planar-memory row contains the
outer factor \(\delta\).  Hence the coupled Volterra map has norm
\(C\delta e^{cS_c}\langle S_c\rangle^M=o(1)\), and a Neumann series gives
the causal inverse with one finite power loss.  Construct a selected
solution first and differentiate its fixed-point equation; the same inverse
then gives the selected first-jet bound.  This proves a solution-jet
statement, not operator differentiability of an inverse on a rough history
space.
\end{proof}

\subsection{The stable sheet at the generated cut}

Retain the scales in
\cref{eq:anchor-main-reserved-scales,eq:anchor-main-generated-cut}.  Thus
\(s_b^\#=\chi_bS_c=c_bS_o\), and put
\begin{equation}
 E_b=e^{(s_b^\#)^2/2},\qquad
 \mathcal A_{b,\delta}=\frac{\delta}{s_b^\#}E_b,\qquad
 \varkappa_\delta=S_o^{-1},\qquad
 \mathfrak b_\delta
 =\varkappa_\delta^2+S_o^{-2}+\sqrt\delta S_o^{7/2}.
 \label{eq:lean-anchor-projective-scale}
\end{equation}
Let \(C_{b,p}\) be the selected repelling history recut at \(\Sigma_b\),
and let \(O_{b,p}\) be the independently generated outer reference there.
Both are histories of the same unmodified anchored equation: the central
cut is a distance \((1-\chi_b)S_c\) from the preparation collar, while the
outer segment has already generated all of its delay values.

\begin{lemma}[Intrinsic and finite sheets on one physical core]
\label[lemma]{lem:lean-anchor-generated-cut-sheets}
There is a compatible chart
\[
 \Phi_{b,p}(\psi_b,u_b),\qquad \Phi_{b,p}(0,0)=O_{b,p},
\]
whose \(u_b\)-column is the oriented event-normal outer unstable column.
Its domain contains the one required radius
\(R_b\mathfrak b_\delta\) ball for all sufficiently small \(\delta\), and
\begin{equation}
 \|\Phi_{b,p}^{-1}(C_{b,p})\|\le C\mathfrak b_\delta,\qquad
 \|D_p^{\rm sel}\{\Phi_{b,p}^{-1}(C_{b,p})\}\|
 \le C\delta^{-M}\mathfrak b_\delta .
 \label{eq:lean-anchor-b-overlap}
\end{equation}
In this chart:

\begin{enumerate}[label=\textup{(\roman*)}]
\item The histories converging to \(E_N^-\) form the intrinsic graph
\begin{equation}
 \mathcal S_b^-(p)
 =\{\Phi_{b,p}(\psi_b,F_{b,p}^-(\psi_b))\},
 \qquad
 \mathfrak m_{b,p}(\Phi_{b,p}(\psi_b,u_b))
 =u_b-F_{b,p}^-(\psi_b).
 \label{eq:lean-anchor-intrinsic-b-sheet}
\end{equation}
The graph is \(C^2\) in state, has a selected first \(p\)-jet, and its
individual \(C^2\) norm is at most \(C\delta^{-M}\).  No small-curvature
claim is made.

\item With the single ordered exponent choice fixed above, a finite graph
\(F_{b,p}^{\rm sel}\) on the same physical core satisfies
\begin{align}
 C_{b,p}&=\Phi_{b,p}
   (\zeta_{C,p},F_{b,p}^{\rm sel}(\zeta_{C,p})),
 \label{eq:lean-anchor-selected-incidence}\\
 \|J_p^1(F_{b,p}^--F_{b,p}^{\rm sel})\|_{C^1}
 &\le C\delta^{Q_*+M_{\rm loss}+4},
 \qquad J_p^1\in\{1,\partial_\nu,D_\eta\}.
 \label{eq:lean-anchor-sheet-comparison}
\end{align}
The finite graph is used only for comparison; the zero set in (i) remains
the intrinsic stable set of the fixed anchored RFDE.

\item With \(k_{b,p}^c\) the full compatible column obtained by
transporting the normalized central \(u\)-column to \(b\),
\begin{equation}
 A_{0b,p}:=D\mathfrak m_{b,p}(C_{b,p})k_{b,p}^c
 \quad\hbox{satisfies}\quad
 c_A\mathcal A_{b,\delta}\le A_{0b,p}
       \le C_A\mathcal A_{b,\delta},
 \qquad |D_p^{\rm sel}\log A_{0b,p}|\le C\delta^{-M}.
 \label{eq:lean-anchor-projective-pairing}
\end{equation}
\end{enumerate}
\end{lemma}

\begin{proof}
The two-scale choice leaves the complete central delay hull strictly inside
the agreement region, while
\cref{eq:anchor-toolkit-generated-hits} makes the outer history literal.
On the central singular trace,
\[
 r=-\frac{\delta s}{2\alpha},\qquad
 w=\frac23-\alpha r^2+\frac{\delta^2}{2\alpha}.
\]
The selected central trace differs from this by its retained graph bound.
On the other hand, \cref{eq:anchor-toolkit-dynamic-overlap} gives the same
\(\delta^2/(2\alpha)\) correction for the exact outer reference.  The
voltage estimate in \cref{prop:anchor-toolkit-reference-hits} and the exact
resource quotient in \cref{prop:anchor-toolkit-resource-phase-quotient}
give, after the event projection and in the triangular quotient chart,
\begin{align*}
 \|\Delta_b^{\rm pl}\|_{b,\sharp}
 &\le C\delta e^{cS_c}\langle S_c\rangle^M,\\
 \|\Delta_b^{\perp}\|_{b,\sharp}+|\Delta_b^a|
 &\le C\mathfrak b_\delta .
\end{align*}
Their selected first \(p\)-derivatives are bounded by
\(C\delta^{-M}\{\delta e^{cS_c}\langle S_c\rangle^M+
\mathfrak b_\delta\}\).
Here \(\Delta_b^a\) is the exact resource-quotient component and the other
two terms are respectively the planar and complementary history
components.  These bounds imply \cref{eq:lean-anchor-b-overlap}; the
quotient denominator is
\(w'_O(-\rho_b)=\delta s_b^\#\{1+o(1)\}\).  Applying the exact event
projection before differentiating and using
\cref{eq:anchor-toolkit-event-jets} proves the two selected first-jet
bounds.  This is an overlap in one compatible chart, not equality of the
two histories.

Use a fixed \(C^2\) clock \(\chi_\beta\) from \(b\) to the regular terminal
section \(L\), equal to the physical fold-time clock on the initial history
window and with its scalar variation supported in a reserved event buffer.
For a path \(Z=(Z_v,Z_w)\), set
\[
 \mathcal H_\beta Z(s)
 =\left(\vartheta\mapsto
 Z_v\!\left[\chi_\beta^{-1}
   \{\chi_\beta(s)+\vartheta\}\right],Z_w(s)\right),
 \qquad
 \mathcal E_p(Z,\beta)
 =Z'-\chi_\beta'\mathscr F^{\rm f}_{N,\delta,p}
       (\mathcal H_\beta Z).
\]
Here \(\mathscr D_{bL}^{0;2}\) consists of paths continuous on the complete
delay enlargement, locally \(W^{1,\infty}\), and \(W^{2,\infty}\) on the
reserved entry, event, and generated terminal windows.  The entry trace is
measured in the weak compatible space \(\mathscr Y_N\), while the terminal
trace lies in the compatible \(C^2\) space \(\mathscr H_L^2\).
\(\mathscr F_{bL}\) is the corresponding distributed-source space with the
action norm of \cref{prop:anchor-toolkit-generated-core}.  For fixed
\(\psi_b\), extend \(\mathfrak m_{L,p}\) to the fixed event tube by the
event-normal projection onto \(L\), constant along the time coordinate; its
restriction and differential on \(L\) are unchanged.  The square residual is
\begin{equation}
 \mathscr R_{bL,p}(Z,\beta,u_b;\psi_b)
 =\left(
 \mathcal E_p(Z,\beta),\
 \mathcal H_\beta Z(0)-\Phi_{b,p}(\psi_b,u_b),\
 r(\mathcal H_\beta Z(T))+r_L,\
 \mathfrak m_{L,p}(\mathcal H_\beta Z(T))
 \right).
 \label{eq:lean-anchor-bL-residual}
\end{equation}
The second row is equality of literal complete histories, not equality of
current values or finite jets.  The typed domain and range of its
linearization are
\begin{equation}
 \mathbb X_{bL}=\mathscr D_{bL}^{0;2}
       \times\mathbb R_\beta\times\mathbb R_u,
 \qquad
 \mathbb Y_{bL}=\mathscr F_{bL}\times\mathscr Y_N
       \times\mathbb R_{\rm ev}\times\mathbb R_{\rm ter}.
 \label{eq:lean-anchor-bL-spaces}
\end{equation}
The quotient dichotomy solves the stable data; the exact resource
shear solves the longitudinal scalar; and the event row removes the time
column.  Denote by
\[
 \mathfrak L^{\rm phys}_{L,p}:
  (y,a,\beta)\longmapsto\widehat y\in\mathscr H_L^2
\]
the inverse resource-quotient shear followed by physical event recutting on
the generated terminal window.  This is a bounded lift by
\cref{prop:anchor-toolkit-resource-phase-quotient}; it includes the current
resource component and the entire terminal voltage history.  If
\(\widehat y_0=\mathfrak L^{\rm phys}_{L,p}(y_0,a[y_0],\beta)\) is the
lift of the particular solution and
\(\widehat h^{u,o}=\mathfrak L^{\rm phys}_{L,p}
(h^{u,o},a[h^{u,o}],0)\) is the lift of the homogeneous unstable column,
the terminal replacement is
\begin{equation}
 \widehat y=\widehat y_0+\widehat h^{u,o}
 \frac{b-D\mathfrak m_{L,p}(O_{L,p})\widehat y_0}
      {D\mathfrak m_{L,p}(O_{L,p})\widehat h^{u,o}}.
 \label{eq:lean-anchor-terminal-replacement}
\end{equation}
Thus the covector and both columns in this formula belong to the same full
terminal-history tangent-dual pair.  The denominator is exactly the scalar
in
\cref{eq:lean-anchor-terminal-crossing}.  The linearization is therefore a
bounded isomorphism, with the sharp \(C\mathscr L_\delta\) coupled loss, in
the action-scaled norms of
\cref{prop:anchor-toolkit-generated-core}.  Its nonlinear two-point bound is
at most
\[
 C_{\rm N}R_b\mathscr L_\delta\mathfrak b_\delta<\frac13,
\]
because
\(\mathscr L_\delta\mathfrak b_\delta=p_o^{-1}+o(1)\) and \(p_o\)
was chosen after the fixed construction radius.  The parameterized
contraction theorem applies because the terminal trace satisfies
\begin{equation}
 \|\operatorname{Tr}_L(X-\bar X)\|
 \le C\delta^{-M}e^{-c\mathscr L_\delta^2}
       R_b\mathfrak b_\delta+C\delta^{-M}e^{-a/\delta}
 <\tfrac1{16}R_{{\rm hl},\delta}.
 \label{eq:lean-anchor-terminal-localization}
\end{equation}
It gives the graph in (i), and literal equality in the
entry row identifies its zeros with forward convergence.  Differentiating
the selected fixed-point equation after it has been solved gives the stated
jets and the polynomial \(C^2\) bound.

Choose a later cut
\(s_c^\#=S_c-\sigma_c\) still inside the exact agreement region and put
\[
 \mathscr A_{bc}=\tfrac12\{(s_c^\#)^2-(s_b^\#)^2\}.
\]
At that cut take the affine terminal row through the selected repelling
history and normalize it on the transported unstable column.  Solving the
same mixed problem gives \(F_{b,p}^{\rm sel}\), and the selected history is
an exact zero, proving \cref{eq:lean-anchor-selected-incidence}.  The two
problems have identical coefficients, history rows, phase row, and stable
coordinate on \([b,c]\); they differ only in the terminal unstable scalar.
Backward propagation of that scalar gives
\[
 \|J_p^1(F_{b,p}^--F_{b,p}^{\rm sel})\|_{C^1}
 \le C\delta^{-M}e^{-\mathscr A_{bc}+cS_c}
       \langle S_c\rangle^M.
\]
Since
\(\mathscr A_{bc}=p_c(1-\chi_b^2)\log(1/\delta)+O(S_c)\), the ordered
inequality for \(p_c\) in \cref{prop:anchor-central-finite-shooting} proves
\cref{eq:lean-anchor-sheet-comparison}.

It remains to compute the typed scalar pairing.  The singular central
column normalized by \(c_{\mathscr H}\mathscr H_\alpha\) is
\[
 n(s)=\binom{I(s)}{E(s)-sI(s)},\qquad
 E(s)=e^{s^2/2},\qquad I(s)=\int_0^sE(t)\,\dd t.
\]
Fixed-\(r\) recutting removes the time component \(I(s)\) and leaves the
transverse component \(E(s)\).  In physical variables its resource
component is \(-\delta^2E(s)\{1+o(1)\}\).  The exact outer resource
quotient divides this by
\(w_O'(-\rho_b)=\delta s_b^\#\{1+o(1)\}\).  If
\(\lambda_{b,p}^{\rm sel}\) is the finite graph row in the full compatible
history dual, the positive-\(\delta\) graph-tangent trace equation and exact
event recut give
\begin{equation}
 \begin{aligned}
  \lambda_{b,p}^{\rm sel}k_{b,p}^c
  &=a_b^0\frac{\delta}{s_b^\#}E_b(1+\mathcal E_{b,p}),
      &0<c_0&\le a_b^0\le C_0,\\
  |\mathcal E_{b,p}|
  &\le C\{\delta e^{cS_c}\langle S_c\rangle^M
               +\mathfrak b_\delta\}=o(1),\\
  |D_p^{\rm sel}\mathcal E_{b,p}|&\le C\delta^{-M}.
 \end{aligned}
 \label{eq:lean-anchor-projective-relative-ledger}
\end{equation}
This is a full row--column pairing; the current covector is applied only
after the resource/history lift.  Replacing the finite row by the intrinsic
row costs
\(C\delta^{Q_*+M_{\rm loss}+4}\|k_{b,p}^c\|\), which is absorbed by the
definition of \(M_{\rm loss}\).
This proves the two-sided estimate and, after differentiating the two
selected trace equations, its logarithmic first-jet bound.
\end{proof}

\begin{lemma}[Action-balanced central intersection]
\label[lemma]{lem:lean-anchor-action-intersection}
Let \(\mathcal M_{0,p}^{\rm int}\) be
\cref{eq:anchor-main-intrinsic-membership}.  After the choices in
\cref{lem:lean-anchor-generated-cut-sheets}, there are radii satisfying
\cref{eq:anchor-main-shooting-radii} for which the action coordinate
\[
 \xi=A_{0b,p}(u-\bar u_p)
\]
obeys
\begin{equation}
 \partial_\xi\mathcal M_{0,p}^{\rm int}\ge\frac12
 \label{eq:lean-anchor-action-orientation}
\end{equation}
on the whole shooting strip.  Its two \(\xi\)-faces have opposite signs,
and its unique zero is the graph in
\cref{eq:anchor-main-intrinsic-root-distance,eq:anchor-main-intrinsic-root-flatness}.
The whole delay hull of this graph remains in the exact central tube.
\end{lemma}

\begin{proof}
At the selected reference,
\cref{eq:lean-anchor-selected-incidence,eq:lean-anchor-sheet-comparison}
give
\begin{equation}
 |J_p^1\mathcal M_{0,p}^{\rm int}(\bar\psi_p,\bar u_p)|
 \le C\delta^{Q_*}.
 \label{eq:lean-anchor-reference-membership}
\end{equation}
Put
\(H_p(\psi,\xi)=\mathcal M_{0,p}^{\rm int}
 (\psi,\bar u_p+A_{0b,p}^{-1}\xi)\).
The graph-tangent estimate in \cref{lem:shared-graph-trace}, the causal
normal inverse in \cref{lem:lean-anchor-central-normal-inverse}, the
polynomial \(C^2\) graph bound, and
\cref{eq:lean-anchor-projective-pairing} give, by the definition of
\(M_{\rm loss}\), on
\[
 \|\psi-\bar\psi_p\|\le c_r\delta^{Q_*+M_{\rm loss}+2},
 \qquad |\xi|\le C_\rho\delta^{Q_*},
\]
one has
\begin{equation}
 \|D_\psi H_p\|\le C\delta^{-M_{\rm loss}},\qquad
 |\partial_\xi H_p-1|\le\frac12,\qquad
 \|D_{(\psi,\xi)}\partial_\xi H_p\|
       \le C\delta^{-M_{\rm loss}}.
 \label{eq:lean-anchor-rectangle-derivatives}
\end{equation}
The physical central displacement produced by \(\xi\) is at most
\((s_b^\#/\delta)|\xi|\).  The graph-tangent and causal normal estimates
bound the contribution of \(\psi-\bar\psi_p\) by
\(C\delta^{-M_{\rm loss}}\|\psi-\bar\psi_p\|\).  With
\(d=V+sU-\alpha U^2\), throughout the
rectangle,
\begin{align}
 \|\Phi_{b,p}^{-1}\mathcal P_{0\to b,p}
       \mathfrak C_{0,p}(\psi,\bar u_p+A_{0b,p}^{-1}\xi)\|
 &\le C\left\{\mathfrak b_\delta
       +\delta^{-M_{\rm loss}}\|\psi-\bar\psi_p\|
       +\frac{s_b^\#}{\delta}|\xi|\right\}
 <R\mathfrak b_\delta,\notag\\
 \sup_{\text{delay hull}}|d|
 &\le C\left\{\delta e^{cS_c}\langle S_c\rangle^M
       +\delta^{-M_{\rm loss}}\|\psi-\bar\psi_p\|
       +\frac{s_b^\#}{\delta}|\xi|\right\}<\delta^\kappa.
 \label{eq:lean-anchor-rectangle-containment}
\end{align}
The ordered inequality \(Q_*>M_{\rm loss}+12\) therefore keeps the event
speed inside its declared margin and puts the \(b\)-hit inside the chart
of \cref{lem:lean-anchor-generated-cut-sheets}.  Choose \(c_r\) small and
then \(C_\rho\) larger than the
constant in \cref{eq:lean-anchor-reference-membership}.  The
\(\psi\)-oscillation is less than one quarter of the face height, while
integration of \cref{eq:lean-anchor-rectangle-derivatives} gives opposite
signs on the two faces.  The intermediate-value theorem gives a zero and
\cref{eq:lean-anchor-action-orientation} gives uniqueness.  Only now do we
apply the parameterized implicit-function theorem.  Dividing
\cref{eq:lean-anchor-reference-membership} by
\(A_{0b,p}\asymp(\delta/s_b^\#)E_b\) proves the distance and first-jet
estimates in the main statement.
\end{proof}

\begin{proof}[Proof of \cref{prop:anchor-central-finite-shooting}]
The common generated section and intrinsic graph are
\cref{lem:lean-anchor-generated-cut-sheets}.  The finite graph contains the
selected repelling history exactly, but is used only to estimate the
intrinsic graph.  The action-balanced intersection and all its radii are
\cref{lem:lean-anchor-action-intersection}.  At a zero the history at
\(\Sigma_b\) lies literally in \(\mathcal S_b^-(p)\); forward uniqueness
and \cref{eq:lean-anchor-intrinsic-b-sheet} give convergence to \(E_N^-\).

Finally, \cref{eq:lean-anchor-positive-selected-flatness}, with its
fixed comparison margin, gives
\[
 \|\psi_A-\bar\psi_p\|
 \le C\delta^{Q_*+M_{\rm loss}+4}.
\]
This is smaller than the algebraic shooting radius, so the actual incoming
history belongs to the graph domain.  No finite terminal row appears in
the definition of that domain or its zero set.
\end{proof}

\begin{proof}[Proof of \cref{thm:anchor-annulus-flat-forgetting}]
The positive history and the intrinsic negative half-line are
\cref{lem:lean-anchor-history-objects}(ii)--(iii).  Forward event holonomy
from \(\Sigma_0\) to the generated cut, together with
\cref{lem:lean-anchor-action-intersection}, gives the graph
\(u=F_{N,\delta,p}^{-,0}(\psi)\).  Differentiating its zero identity gives
the pullback formula
\cref{eq:anchor-central-conormal-forward-pullback}; only a forward
derivative occurs.

For a finite tail selector, reduce the infinite and finite problems to the
same core before the buffer.  Their fixed-point maps differ only by the
future unstable integral and the fading-history tail.  The capped action in
\cref{lem:lean-anchor-history-objects}(iii) is at least
\(c\log(1/\delta)/\delta\).  The mean-value fixed-point identity, followed
by one state or selected parameter derivative, gives
\cref{eq:anchor-flat-envelope,eq:anchor-flat-jet-list}.  Pull both graph rows
back to the same entry space and normalize them on the same unstable column;
the quotient rule gives
\cref{eq:anchor-flat-conormal-comparison}.  This proves the theorem.
\end{proof}

\subsection{The physical gap as a direct perturbation of the selected gap}

Fix one structural-ball admissible preparation, solely as a comparison
device; by the convention at the start of this section it is precisely
\(\mathbf p_{\rm sel}\), not a new choice.  Let
\(A^{\rm sel}_{0,p}\) and \(B^{\rm sel}_{0,p}\) be its attracting and
repelling selected histories on \(\Sigma_0\).  In the tubular chart write
\[
 A^{\rm sel}_{0,p}=\mathfrak C_{0,p}(\bar\psi_p,u^a_p),
 \qquad
 B^{\rm sel}_{0,p}=\mathfrak C_{0,p}(\bar\psi_p,u^r_p).
\]
Both traces belong to the same exact invariant-history lift.  The normal
coordinate of that lift is zero before parameter trivialization and hence
is the same \(\bar\psi_p\) in the displayed chart.  By
\cref{eq:anchor-central-chart-identities},
\begin{equation}
 \bar u_p=u_p^r,\qquad
 \mathscr D_N^{\rm sel}=u^a_p-u^r_p.
 \label{eq:lean-anchor-selected-gap-coordinate}
\end{equation}

\begin{lemma}[Selected-to-intrinsic gap comparison]
\label[lemma]{lem:lean-anchor-gap-comparison}
Uniformly in \(N\) and on the common \((\nu,\eta)\)-box, the single ordered
choice in \cref{prop:anchor-central-finite-shooting} gives
\begin{equation}
 \bigl|J_p^1\{G^{\rm anc}_{N,\delta}
          -\mathscr D_N^{\rm sel}\}\bigr|
 \le C\delta^{Q_*},
 \qquad J_p^1\in\{1,\partial_\nu,D_\eta\}.
 \label{eq:lean-anchor-gap-C1-comparison}
\end{equation}
The left member is the gap of the fixed anchored RFDE; the preparation
occurs only in the comparison term on the right.
\end{lemma}

\begin{proof}
Write \(A^0_p=\mathfrak C_{0,p}(\psi_A,u_A)\).  The positive common-core
comparison \cref{eq:lean-anchor-positive-selected-flatness} gives
\begin{equation}
 \bigl\|J_p^1\{(\psi_A,u_A)-(\bar\psi_p,u^a_p)\}\bigr\|
 \le C\delta^{Q_*+M_{\rm loss}+4}.
 \label{eq:lean-anchor-positive-selected-comparison}
\end{equation}
Here and below a superalgebraic capped-action estimate has been weakened to
the prescribed finite order.

On the negative side, the finite graph in
\cref{lem:lean-anchor-generated-cut-sheets} contains the selected repelling
history.  Pull both finite and intrinsic graphs to \(\Sigma_0\) by their
common forward membership composition.  No history holonomy is inverted.
The two defining functions differ by the right side of
\cref{eq:lean-anchor-sheet-comparison}, times only a finite central trace
loss, while their \(u\)-derivatives are bounded below by
\(c(\delta/s_b^\#)E_b\).  Choose the comparison order larger than this
finite loss and the reciprocal action exponent.  The scalar implicit
estimate then gives
 \begin{equation}
 \bigl|J_p^1\{F_{N,\delta,p}^{-,0}(\bar\psi_p)-u^r_p\}\bigr|
 \le C\delta^{Q_*+4}.
 \label{eq:lean-anchor-negative-selected-comparison}
 \end{equation}
The individual graph may have polynomial slope.  The margin
\(M_{\rm loss}+4\) in
\cref{eq:lean-anchor-positive-selected-comparison} ensures that this
polynomial loss times \(\|\psi_A-\bar\psi_p\|\) is still
\(O(\delta^{Q_*})\).  Then
\begin{align*}
 G^{\rm anc}-\mathscr D_N^{\rm sel}
  &=(u_A-u^a_p)
    -\{F_{N,\delta,p}^{-,0}(\psi_A)-u^r_p\}
\end{align*}
by \cref{eq:anchor-gap-coordinate-transversality,eq:lean-anchor-selected-gap-coordinate}.
The preceding two estimates and the mean-value theorem prove the value
bound.  Differentiate the three already constructed selected families and
use their mixed state--parameter bounds to obtain the two first-parameter
bounds.  No derivative of a rough-history inverse is taken.
\end{proof}

\begin{proof}[Proof of \cref{prop:anchor-gap-comparison}]
The chart identities give the exact formula
\[
 G^{\rm anc}_{N,\delta}
 =u_A-F_{N,\delta,p}^{-,0}(\psi_A).
\]
Thus \(G^{\rm anc}=0\) if and only if the incoming history equals the
point of the intrinsic stable sheet with the same full normal/history
coordinate.  Its past is the inward branch of \(W^u(E_N^+)\), its future
converges to \(E_N^-\), and RFDE forward uniqueness joins them into one
orbit.  The converse follows by recutting any such orbit at \(\Sigma_0\).
This proves the exact zero-fiber statement and the declared local
uniqueness.

For the estimates, use the fixed \(Q_*>M_{\rm loss}+12\).  The
structural-ball version of the finite-gap bridge is
\cref{eq:full-ball-gap-response}; together with
\cref{eq:gap-bridge-expansion,eq:gap-bridge-nu} it gives
\begin{align*}
 \frac{\mathscr D_N^{\rm sel}}\delta
 &=\sqrt{2\pi}(\nu-\nu_{0,N})+O(\delta),\\
 \partial_\nu\mathscr D_N^{\rm sel}
 &=\delta\sqrt{2\pi}+O(\delta^2),\\
 \|D_\eta\mathscr D_N^{\rm sel}
       +\delta^2\sqrt{2\pi}\Lambda_N\|
 &\le C\{\delta^3+\delta^2\|\eta\|\}.
\end{align*}
Insert \cref{eq:lean-anchor-gap-C1-comparison}.  This proves
\cref{eq:anchor-physical-gap-estimates}.  The ordered choice in
\cref{prop:anchor-central-finite-shooting} makes both the Gaussian
selected-gap tail and the common-core action dominate \(Q_*\) and all finite
first-jet losses.  The fixed-annulus comparison has the separate
capped-action rate.  No second preparation exponent or endpoint-tail
recalculation is used.
\end{proof}

\begin{proof}[Proof of \cref{thm:anchored-physical-root-conormal}]
The first two estimates in \cref{eq:anchor-physical-gap-estimates} give
opposite signs at the ends of a common \(\nu\)-window and
\[
 \partial_\nu G^{\rm anc}_{N,\delta}
 \ge\tfrac12\delta\sqrt{2\pi}>0.
\]
Hence there is exactly one root in that window.  By
\cref{prop:anchor-gap-comparison}, it is the complete-history heteroclinic
root of the fixed anchored RFDE and is independent of the preparation used
in \cref{lem:lean-anchor-gap-comparison}.

Implicit differentiation yields
\[
 D_\eta\nu^{\rm anc}_{c,N}
 =-\frac{D_\eta G^{\rm anc}}{\partial_\nu G^{\rm anc}}
 =\delta\Lambda_N+O(\delta^2+\delta\|\eta\|)
\]
in the full dual norm.  Multiplication by \(\delta^2\), followed by
integration on \(t\mapsto t\eta\), proves
\cref{eq:anchor-physical-root-jets,eq:anchor-physical-root-increment}.
After centering and dividing by \(\delta^3\), this is
\cref{eq:anchor-physical-conormal-limit}.

The tracking assertions are already part of the two history objects in
\cref{lem:lean-anchor-history-objects}; the central estimate follows from
\cref{lem:shared-graph-trace,lem:lean-anchor-central-normal-inverse}.
At \(b\), the root history and the negative chart center lie on the same
stable graph and differ by \(O(\mathfrak b_\delta)\).  The generated-core
stable dichotomy contracts this difference over the \(b\)-to-\(c\) action:
\[
 C\mathfrak b_\delta e^{-c\mathscr L_\delta^2}=O(\delta^P)
 \quad\text{for every fixed }P.
\]
It may therefore be joined to the outer chart-center bound on every fixed
repelling subannulus.  On a fixed outer subannulus the
resource speed is \(O(\delta^2)\) in original time, so its traversal takes
at least \(c\delta^{-2}\).  At the root, forward uniqueness identifies the
two one-sided restrictions with the same orbit.  Finally, the factorization,
right inverse, one-delay no-go, and nonvanishing examples concern the same
covector \(\Lambda_N\); they therefore apply to the limiting intrinsic
conormal.
\end{proof}

\begin{proof}[Proof of
  \cref{prop:anchor-physical-naturality-composition}]
Two regular sections reached by forward holonomy describe the same local
regular parameter zero set.  On that set their differentials have the same
tangent kernel and hence span the same conormal line.  If a passage is split
into forward maps
\(P_j\circ\cdots\circ P_1\), the ordinary chain rule gives
\(D(P_j\circ\cdots\circ P_1)=DP_j\cdots DP_1\).  No inverse history
holonomy occurs.  A nonsingular parameter change acts by cotangent pullback.

For each admitted anchor multiplier the preceding root calculation gives a
weighted conormal within \(C\delta\) of \((1,-\Lambda_N)\), after centering
at that multiplier's own baseline.  Uniformity over the bounded anchor
class follows from the common constants in
\cref{lem:lean-anchor-history-objects}.  The triangle
inequality gives the asserted \(2C\delta\) comparison, without identifying
the two finite-\(\delta\) roots.
\end{proof}
